\begin{examplebox}
	\textbf{\textcolor{CExample}{例 1（基础）}}

	\textbf{\textcolor{CExample}{题目：}}
	在$\triangle ABC$中，$\angle A = 40^\circ$，$\angle B = 60^\circ$；在$\triangle DEF$中，$\angle D = 40^\circ$，$\angle E = 60^\circ$。求证：$\triangle ABC\sim\triangle DEF$。

	\textbf{\textcolor{CExample}{解题步骤：}}
	\begin{description}[style=nextline, leftmargin=2.8em, labelsep=0.5em, itemsep=0.45em]
		\item[\textbf{第一步（找等角）：}] 由已知，$\angle A = \angle D = 40^\circ$，$\angle B = \angle E = 60^\circ$。
			\par\textbf{理由：} 题目条件直接给出。
		\item[\textbf{第二步（应用判定）：}] 因为两个角对应相等，由AA（角角）判定定理，可以判定两三角形相似，无需再验证第三个角。
			\par\textbf{理由：} 相似三角形判定AA。
		\item[\textbf{第三步（写出结论）：}] 所以 $\triangle ABC \sim \triangle DEF$。
			\par\textbf{理由：} 满足AA判定条件。
	\end{description}

	\textbf{\textcolor{CExample}{关键结论：}}
	$\triangle ABC \sim \triangle DEF$。

	\medskip
	\textbf{\textcolor{CExample}{例 2（稍变形）}}

	\textbf{\textcolor{CExample}{题目：}}
	在$\triangle ABC$中，$AB=3$，$BC=4$，$\angle B=50^\circ$；在$\triangle DEF$中，$DE=6$，$EF=8$，$\angle E=50^\circ$。求证：$\triangle ABC\sim\triangle DEF$。

	\textbf{\textcolor{CExample}{解题步骤：}}
	\begin{description}[style=nextline, leftmargin=2.8em, labelsep=0.5em, itemsep=0.45em]
		\item[\textbf{第一步（算边比）：}] 因为 $AB=3$, $DE=6$，所以 $\frac{AB}{DE} = \frac12$；同理 $\frac{BC}{EF} = \frac12$，故 $\frac{AB}{DE} = \frac{BC}{EF}$。
			\par\textbf{理由：} 代入长度数值。
		\item[\textbf{第二步（找夹等角）：}] 已知 $\angle B = \angle E = 50^\circ$，且$\angle B$是$AB$与$BC$的夹角，$\angle E$是$DE$与$EF$的夹角，满足夹角条件。
			\par\textbf{理由：} 题目条件。
		\item[\textbf{第三步（SAS判定）：}] 两边对应成比例且夹角相等，由SAS判定得$\triangle ABC \sim \triangle DEF$。
			\par\textbf{理由：} 相似三角形判定SAS。
	\end{description}

	\textbf{\textcolor{CExample}{关键结论：}}
	$\triangle ABC \sim \triangle DEF$。

	\medskip
	\textbf{\textcolor{CExample}{例 3（综合应用）}}

	\textbf{\textcolor{CExample}{题目：}}
	如图（略），在$\triangle ABC$中，$DE\parallel BC$，$D$在$AB$上，$E$在$AC$上。若$AD=2$，$DB=3$，$DE=4$，求$BC$的长。

	\textbf{\textcolor{CExample}{解题步骤：}}
	\begin{description}[style=nextline, leftmargin=2.8em, labelsep=0.5em, itemsep=0.45em]
		\item[\textbf{第一步（证相似）：}] 因为$DE\parallel BC$，由平行截判定，得$\triangle ADE\sim\triangle ABC$。
			\par\textbf{理由：} 平行于三角形一边的直线截得的三角形与原三角形相似。
		\item[\textbf{第二步（列比例式）：}] 由相似，对应边成比例：$\frac{AD}{AB} = \frac{DE}{BC}$。又$AB = AD+DB = 2+3=5$。
			\par\textbf{理由：} 相似三角形对应边成比例。
		\item[\textbf{第三步（解方程）：}] 代入数值：$\frac{2}{5} = \frac{4}{BC}$，解得$BC=10$。
			\par\textbf{理由：} 比例方程求解。
	\end{description}

	\textbf{\textcolor{CExample}{关键结论：}}
	$BC=10$。
\end{examplebox}
